\chapter{Marco teórico}
\label{ch:marco-teorico}

Antes de adentrarnos en el estudio de los ceros de los polinomios ortogonales de Sobolev, es esencial establecer el marco teórico que usaremos a lo largo del trabajo. En este sentido, comenzaremos por definir conceptos necesarios a continuación.

\section{Polinomios Ortogonales}

\begin{definition}\label{def:espacio-polinomios}
  Sea $\mathbb{K}$ un cuerpo (normalmente $\mathbb{R}$ o $\mathbb{C}$).
  El \emph{espacio de polinomios en una variable} se define como
  \[
    \mathbb{P}[z]\;=\;
    \Bigl\{\,p(z)=\sum_{k=0}^{n} a_k\,z^{\,k}\;:\;
    a_k\in\mathbb{K},\;n\in\mathbb{N}\Bigr\}.
  \]
  Para cada $n\in\mathbb{N}$ denotamos
  $
    \mathbb{P}_{n}[z]=\{p\in\mathbb{P}[z]:\deg p\le n\},
  $
  de modo que $\mathbb{P}[z]=\bigcup_{n=0}^{\infty}\mathbb{P}_{n}[z]$.
  El conjunto $\{1,z,z^{2},\dots,z^{n}\}$ forma una base de
  $\mathbb{P}_{n}[z]$, luego $\dim\!\bigl(\mathbb{P}_{n}[z]\bigr)=n+1$.
\end{definition}

\begin{definition}\label{def:producto-interno-general-medida}
  Sea $\mu$ una medida de Borel positiva con soporte infinito $\operatorname{supp}(\mu)=\{z\;\mu(B_r(z))>0\;\forall r>0\}$. Definimos en $\mathbb{P}[z]$ un \emph{producto interno} cualquiera por
  \[
    \langle p,q\rangle
    \;=\;
    \int p(z)\,\overline{q(z)}\,\mathrm d\mu(z),
    \qquad p,q\in\mathbb{P}[z].
  \]
  Esta forma sesqui-lineal satisface:
  \begin{enumerate}
    \item $\langle p,q\rangle=\overline{\langle q,p\rangle}$,
    \item $\langle \alpha p+\beta r,q\rangle
            =\alpha\langle p,q\rangle+\beta\langle r,q\rangle$,
    \item $\langle p,p\rangle>0$ si $p\neq0$.
  \end{enumerate}
\end{definition}

\begin{definition}\label{def:prehilbert_space}
  Un \textbf{espacio prehilbertiano} o \textbf{espacio con producto interno} $(X, \langle \cdot, \cdot \rangle)$ es un espacio vectorial $X$ dotado de un producto interno $\langle \cdot, \cdot \rangle: X \times X \rightarrow \mathbb{C}$ (o $\mathbb{R}$), que satisface las siguientes propiedades para todo $x, y, z \in X$ y $\alpha \in \mathbb{C}$ (o $\mathbb{R}$):
  \begin{enumerate}
    \item $\langle x, y \rangle = \overline{\langle y, x \rangle}$ (simetría hermitiana)
    \item $\langle \alpha x + \beta y, z \rangle = \alpha \langle x, z \rangle + \beta \langle y, z \rangle$ (linealidad en la primera variable)
    \item $\langle x, x \rangle \geq 0$ y $\langle x, x \rangle = 0$ si y solo si $x = 0$ (definida positiva)
  \end{enumerate}

  Un \textbf{espacio de Hilbert} es un espacio prehilbertiano que además es completo respecto a la norma inducida por el producto interno $\|x\| = \sqrt{\langle x, x \rangle}$, es decir, toda sucesión de Cauchy converge a un elemento del espacio.
\end{definition}

\begin{definition}\label{def:gram-schmidt}
  Sea $\{v_0,v_1,\dots,v_n\}$ un sistema generador de un
  subespacio $W\subset V$ con producto interno.
  Construimos recursivamente
  \[
    u_0=v_0,\qquad
    u_k= v_k-\sum_{j=0}^{k-1}
    \frac{\langle v_k,u_j\rangle}{\langle u_j,u_j\rangle}\,u_j,
    \;k=1,\dots,n.
  \]
  El conjunto ortonormal
  $\displaystyle e_k=\frac{u_k}{\|u_k\|}$ satisface
  $\langle e_i,e_j\rangle=\delta_{ij}$ y
  $\operatorname{span}\{e_0,\dots,e_n\}=W$.
\end{definition}

El proceso de Gram-Schmidt es fundamental para construir bases ortonormales, una herramienta esencial en el estudio de los polinomios ortogonales.

Dado un producto interno $\langle\cdot,\cdot\rangle$ en $\mathbb{P}[z]$,
la \emph{sucesión de polinomios ortonormales}
$\{\varphi_n\}_{n\ge 0}$ se obtiene aplicando
el algoritmo de Gram-Schmidt a la base monomial $\{1,z,z^2,\dots\}$:
\[
  \deg \varphi_n=n,\qquad
  \langle \varphi_n,\varphi_k\rangle=\delta_{nk}.
\]
Cada $P_n$ es único salvo su signo si se exige que su coeficiente principal sea positivo.


\begin{definition}\label{def:ceros}
  Sea $p\in\mathbb{P}[z]\setminus\{0\}$.
  El \emph{conjunto de ceros} de $p$ es
  \[
    Z(p)=\bigl\{\,z\in\mathbb{C}\;:\;p(z)=0\bigr\}.
  \]
  Además, para la sucesión de polinomios ortonormales \(\{\varphi_n\}_{n\ge0}\) asociada al producto interno, definimos el conjunto de ceros $Z$
  \[
    Z \;=\; \bigcup_{n=0}^{\infty} Z_n
    \;=\;\bigl\{\,z\in\mathbb{C} : \exists\,n\in\mathbb{N},\;P_n(z)=0\bigr\}.
  \]
  donde $Z_n = Z(P_n)$ es el conjunto de ceros del polinomio $P_n$.
\end{definition}

La localización y propiedades de estos conjuntos de ceros son un tema central en la teoría de polinomios ortogonales. Para facilitar su estudio, a menudo se recurre a representaciones matriciales.

\section{Matrices de momentos asociadas a productos internos}

\begin{definition}\label{def:hermitian_matrix}
  \cite{Horn_Johnson_2012} Una matriz $A \in \mathcal{M}_{n}(\mathbb{C})$ se denomina \textbf{hermitiana} si $A = A^*$, donde $A^*$ denota la matriz conjugada transpuesta de $A$. En el caso real, una matriz simétrica satisface $A = A^t$.

  Una matriz hermitiana $A$ se dice:
  \begin{enumerate}
    \item \textbf{Definida positiva} (HPD) si para todo vector no nulo $x \in \mathbb{C}^n$, se cumple $xAx^* > 0$. Equivalentemente, si todos sus autovalores son estrictamente positivos.

    \item \textbf{Semidefinida positiva} (HPSD) si para todo vector $x \in \mathbb{C}^n$, se cumple $xAx^* \geq 0$. Equivalentemente, si todos sus autovalores son no negativos.
  \end{enumerate}

  Para una matriz hermitiana, el producto escalar $\langle x, y \rangle_A = xAy^*$ define un producto interno si y solo si $A$ es definida positiva.
\end{definition}

\begin{definition}\label{def:moment_matrix}
  Sea \(\langle\!\cdot,\cdot\!\rangle\) un producto interno en el espacio de los polinomios
  \(\mathbb{P}[z]\) y considérese la base algebraica
  \(\mathcal{B}=\{\,z^{k}\,\}_{k\ge 0}\).
  La \emph{matriz de momentos} asociada a \(\langle\!\cdot,\cdot\!\rangle\) es la matriz
  infinita
  \[
    M=\bigl(c_{i,j}\bigr)_{i,j\ge 0},
    \qquad
    c_{i,j}:=\bigl\langle z^{i},\,z^{j}\bigr\rangle ,\; i,j\in\mathbb N_{0}.
  \]

  En particular, si existe una medida de Borel \(\mu\) con soporte infinito tal que
  \(
  \langle p,q\rangle=\displaystyle\int p(z)\,\overline{q(z)}\,\mathrm d\mu(z),
  \)
  entonces
  \[
    c_{i,j}=\int z^{\,i}\,\overline{z^{\,j}}\,\mathrm d\mu(z)\, .
  \]%
\end{definition}

Esta matriz de momentos encapsula la información del producto interno respecto a la base monomial.

Desde una perspectiva matricial, que facilita el análisis algebraico, podemos considerar que una matriz infinita HPD $M$ induce en $\mathbb{P}[z]$ el siguiente producto:
\begin{equation}
  \langle p(z), q(z) \rangle_M = v M w^{*},
\end{equation}
donde $v = (v_0, v_1, \ldots, v_n, 0, \ldots)$, $w = (w_0, w_1, \ldots, w_m, 0, \ldots) \in c_{00} = \{(x_n)_{n=1}^{\infty} \in \mathbb{C}^{\mathbb{N}} : \exists n_0 \in \mathbb{N} \quad | \quad x_n=0 ,\forall n\geq n_0\}$ para polinomios $p(z) = \sum_{k=0}^n v_k z^k$ y $q(z) = \sum_{k=0}^m w_k z^k$. La matriz de momentos $M = (c_{ij})^{\infty}_{i,j=0}$ definida por: $c_{ij} = \langle z^i, z^j \rangle_{M}$, induce una norma Hilbertiana denotada por: $\|p(z)\|^{2}_{M} = \langle p(z), p(z)\rangle_{M},\forall p(z) \in P[z]$

%-------------------------------------------------
\begin{definition}\label{def:section_matrix}
  Sea \(A=\bigl(a_{ij}\bigr)_{\,i,j\ge 0}\) una matriz infinita.
  Para cada \(n\in\mathbb{N}\) definimos su \emph{sección principal de orden \(n\)} (también llamada
  \emph{truncación} o \emph{submatriz principal}) como

  \[
    A_n \;=\;
    \bigl(a_{ij}\bigr)_{\,0\le i,j\le n-1}\;=\;
    \begin{pmatrix}
      a_{00}     & a_{01}     & \dots  & a_{0n-1}       \\
      a_{10}     & a_{11}     & \dots  & a_{1n-1}       \\
      \vdots     & \vdots     & \ddots & \vdots         \\
      a_{(n-1)0} & a_{(n-1)1} & \dots  & a_{(n-1)(n-1)}
    \end{pmatrix}.
  \]

  Estas secciones permiten “aproximar” \(A\) por operadores de rango finito y realizar cálculos numéricos.
\end{definition}
%-------------------------------------------------

Las secciones principales son cruciales para el análisis numérico y la aproximación de operadores definidos por matrices infinitas.

\section{Operador multiplicación por z}

\begin{definition}\label{def:d_application}
  El operador multiplicación por $z$, lo denotaremos como $\mathcal{D}$, y es la aplicación $\mathcal{D}:\mathbb{P}[z] \mapsto \mathbb{P}[z]$ tal que: $D(p(z))=zp(z)$.
  Esta aplicación es lineal ya que
  $$\mathcal{D}(\alpha p(z) + \beta q(z))=z(\alpha p(z) + \beta q(z))=\alpha \mathcal{D}(p(z)) + \beta \mathcal{D}(q(z))$$
  \noindent Denotaremos como $\mathbf{D}$ la representación matricial del operador multiplicación por $z$ en la base de los polinomios ortonormales.
\end{definition}

\begin{definition}\label{def:bounded_operator}
  El operador multiplicación $\mathcal{D}$ se dice \textbf{acotado} en $(\mathbb{P}[z], \langle \cdot, \cdot \rangle)$ si existe una constante $C > 0$ tal que
  \[
    \|\mathcal{D}(p)\| \leq C \|p\|, \qquad \forall p \in \mathbb{P}[z].
  \]
  La menor constante $C$ que satisface esta desigualdad se denomina \textbf{norma del operador} y se denota por $\|\mathcal{D}\|$.
\end{definition}

La acotación de este operador tiene importantes consecuencias para la localización de los ceros de los polinomios ortogonales. A continuación, se presenta cómo se representa este operador en la base de polinomios ortonormales.

\begin{definition}\label{def:d_representation}
  Sea \(\{\varphi_n\}_{n\ge0}\) una base ortonormal de \(\mathbb{P}[z]\) respecto a un
  producto interior \(\langle\!\cdot,\cdot\!\rangle\).
  $z\varphi_n(z)$ es un polinomio de grado $n+1$ que se puede expresar en términos de la base ortonormal como:
  $$
    z\varphi_n(z)= d_{0,0}\varphi_0(z)+d_{1,0}\varphi_1(z)+\dots+d_{n,0}\varphi_n(z)+d_{n,n+1}\varphi_{n+1}(z)
  $$
  \noindent Esta expresión recurrente define la matriz infinita $\mathbf{D}=(d_{i,j})_{i,j=0}^{\infty}$ que representa al operador multiplicación $\mathcal{D}$ respecto de la base de polinomios ortonormales. Denotamos por $\mathbf{D}_n$ a la sección de orden $n$.
\end{definition}

\begin{proposition}\label{prop:matriz_hessenberg}
  Para todos los índices \(i,j\ge0\) se verifica
  \[
    d_{i,j}=0
    \quad\text{si }\,i>j+1.
  \]
  En consecuencia
  \(\mathbf{D}\) es una matriz Hessenberg superior infinita:

  \[
    \mathbf{D}=\begin{pmatrix}
      d_{0,0} & d_{0,1} & d_{0,2} & d_{0,3} & \cdots \\
      d_{1,0} & d_{1,1} & d_{1,2} & d_{1,3} & \cdots \\
      0       & d_{2,1} & d_{2,2} & d_{2,3} & \cdots \\
      0       & 0       & d_{3,2} & d_{3,3} & \cdots \\
      \vdots  & \vdots  & \vdots  & \vdots  & \ddots
    \end{pmatrix}.
  \]
  El patrón de ceros proviene del hecho de que
  \(\deg(z\varphi_j)=j+1\), de modo que al expandir \(z\varphi_j\) en la base
  solo aparecen \(\varphi_0,\dots,\varphi_{j+1}\).
\end{proposition}

La estructura de Hessenberg de la matriz $\mathbf{D}$ es una propiedad clave que simplifica muchos análisis. Esta matriz posee varias propiedades importantes, especialmente en relación con los ceros de los polinomios ortonormales.

\paragraph{Propiedades de la matriz $\mathbf{D}$}
\begin{itemize}
  \item Si el producto interior procede de una medida real en un intervalo
        (caso clásico), la relación de tres términos
        \(z\varphi_j=a_{j+1}\varphi_{j+1}+b_j\varphi_j+a_j\varphi_{j-1}\)
        hace que \(\mathbf{D}\) sea \emph{tridiagonal y simétrica} (matriz de Jacobi),
        un caso particular dentro de las Hessenberg.
  \item Las secciones principales
        \(\mathbf{D}_n=(d_{i,j})_{0\le i,j\le n}\) son matrices \((n+1)\times(n+1)\)
        cuya transpuesta tiene como autovalores los ceros del polinomio
        \(\varphi_{n+1}\).
        Esto conecta la teoría espectral de \(\mathbf{D}\) con la localización de ceros.
  \item El determinante de $\mathbf{D}_n -z \mathbf{I}_n$ denotado por $\vert \mathbf{D}_n -z \mathbf{I}_n\vert $ donde $\mathbf{I}_n$ es la matriz identidad de dimensiones $(n+1)\times (n+1)$ es un polinomio proporcional a $\varphi_{n+1}(z)$, es decir, tienen las misma raíces.
\end{itemize}

La conexión entre los autovalores de las secciones de $\mathbf{D}$ y los ceros de los polinomios ortonormales es un resultado fundamental, cuya demostración se detalla a continuación.

\begin{proposition}\label{prop:eigenvalues_zeros}
  Los ceros de los polinomios ortogonales $\varphi_{n+1}(z)$ son exactamente los autovalores de la matriz transpuesta $\mathbf{D}_n^t$, es decir,
  \[
    \sigma(\mathbf{D}_n^{\,t}) = Z(\varphi_{n+1}) = \{z \in \mathbb{C} : \varphi_{n+1}(z) = 0\}.
  \]
\end{proposition}


\begin{proof}
  \noindent Sea $z_0$ una raíz de $\varphi_{n+1}(z)$, es decir, $\varphi_{n+1}(z_0)=0$. Entonces:

  $$
    z_0\varphi_0(z_0)= d_{0,0}\varphi_0(z_0)+d_{1,0}\varphi_1(z_0)
  $$
  $$
    z_0\varphi_1(z_0)= d_{0,1}\varphi_0(z_0)+d_{1,1}\varphi_1(z_0)+d_{2,1}\varphi_{2}(z_0)
  $$
  $$
    ----------------
  $$
  $$
    z_0\varphi_n(z_0)= d_{0,n}\varphi_0(z_0)+d_{1,n}\varphi_1(z_0)+\dots+d_{n,n}\varphi_n(z_0)+d_{n+1,n}\varphi_{n+1}(z_0)
  $$

  \noindent Puesto que $\varphi_{n+1}(z_0)=0$, la última igualdad queda:
  $$
    z_0\varphi_n(z_0)= d_{0,0}\varphi_0(z_0)+d_{1,0}\varphi_1(z_0)+\dots+d_{n,0}\varphi_n(z_0).
  $$

  \noindent Esto se puede expresar matricialmente de la siguiente forma
  $$
    z_0\begin{pmatrix} \varphi_0(z_0) \\ \varphi_1(z_0) \\  \vdots \\ \varphi_n(z_0) \end{pmatrix} = \mathbf{D}^{t}_n \begin{pmatrix} \varphi_0(z_0) \\ \varphi_1(z_0) \\  \vdots \\ \varphi_n(z_0) \end{pmatrix}
  $$

  \noindent lo cual significa que $z_0$ es autovalor de la matriz $\mathbf{D}^{t}_n$ siendo el autovector \newline
  $(\varphi_0(z_0),\varphi_1(z_0),\dots,\varphi_n(z_0))$ (que es no nulo ya que $\varphi_0(z)$ es constante).

  \noindent Por lo tanto, todos las raíces de $\varphi_n$ son autovalores de la matriz $\mathbf{D}^{t}_n$, y por tanto de $\mathbf{D}_n$.
\end{proof}

Este resultado establece una correspondencia directa entre el problema algebraico de encontrar autovalores y el problema analítico de encontrar ceros de polinomios.

\noindent Como consecuencia, el determinante $|\mathbf{D}_n -z \mathbf{I}_n|$ es un polinomio ortogonal mónico.

El enfoque matricial no solo se limita a la representación del operador multiplicación, sino que también permite caracterizar los propios polinomios ortogonales.

Siguiendo con el enfoque matricial propuesto por C. Escribano y R. Gonzalo, sabemos que los polinomios mónicos ortogonales de Sobolev se pueden obtener a partir de la matriz del operador multiplicación $\mathcal{D}$ y la matriz de momentos asociada al producto interno inducido por $\mathbf{M}$.

\begin{definition}\label{def:spectral_radius}
  \cite{dunford1963linear} El \textbf{radio espectral} de una matriz cuadrada $A \in \mathcal{M}_{n}(\mathbb{C})$, denotado por $\rho(A)$, se define como el máximo de los valores absolutos de sus autovalores:
  \begin{equation}
    \rho(A) = \max\{|\lambda| : \lambda \text{ es autovalor de } A\}
  \end{equation}
\end{definition}

\begin{definition}\label{def:singular_values}
  Sea $A \in \mathcal{M}_{n}(\mathbb{C})$ matriz finita, la matriz $\mathbf{A}^{*}\mathbf{A}$  ( $\mathbf{A}^{t}\mathbf{A}$ en el caso de matrices con entradas reales) es una matriz hermitiana de orden n y semidefinida positiva. Por lo tanto, tiene $n$ autovalores mayores o iguales que cero contados con su multiplicidad
  $$
    0 \leq \lambda_{n}\leq \dots \leq \lambda_{2} \leq \lambda_{1}
  $$


  \noindent Los valores singulares de $\mathbf{A}$ son:
  $$
    \sigma_{k}= \sqrt{\lambda_{k}}
  $$
  \noindent siendo:
  $$
    0 \leq \sigma_{n} \leq \dots \sigma_{2} \leq \sigma_{1}
  $$

  \noindent Utilizando el criterio de Rayleigh y el Principio de min-max de  Courant-Fischer \cite{gallier2017spectral} se tiene que estos valores singulares tienen relación con la norma de aplicaciones lineales y continuas cuando consideramos la norma euclídea en los espacios $\mathbb{C}^n$. De hecho,
  \begin{enumerate}
    \item $$
            \sigma_{1}^2 =
            \max \{ <\mathbf{A}^{*	}\mathbf{A}x,x>: \Vert x \Vert =1 \}
            = \max \{\Vert \mathbf{A} x \Vert^2 : \Vert x \Vert =1\}=
            \Vert \mathbf{A} \Vert^2
          $$
    \item Principio min-max: el segundo mayor valor singular de una matriz es el mínimo de las normas de las restricciones de la aplicación $\mathbf{A}$ a los espacios de codimnesión $1$.
          $$
            \sigma_{2}^2 = \min_{V: cod(V)=1} \max \{\Vert \mathbf{A}x \Vert^2 , x\in V, \Vert x \Vert =1\}= \min_{V: cod(V)=1} \Vert \mathbf{A}/V \Vert^2
          $$

          \noindent siendo $\mathbf{A}/V$ la restricción de la aplicación lineal $\mathbf{A}$ al subespacio $V$.

  \end{enumerate}
\end{definition}

Los valores singulares, especialmente el mayor, están directamente relacionados con la norma del operador, lo que es crucial para el análisis de la acotación. Para simplificar la discusión futura, introducimos la siguiente notación.

\begin{definition}\label{def:notation_matrix_sections}
  Sea $\mathbf{D}_{n}$ la sección de la matriz infinita $\mathbf{D}$ que representa al operador $\mathcal{D}$ multiplicación por $z$ en una base de polinomios ortogonales asociada a un producto internodefinido en el espacio de polinomios $\mathbb{P}[z]$. Definimos:
  \begin{enumerate}
    \item Los valores singulares de $\mathbf{D}_{n}$ se denotan como $\sigma_{1,n} \ge \sigma_{2,n} \ge \ldots, \sigma_{n,n}$
    \item El radio espectral como $\rho(\mathbf{D}_{n})$
    \item El conjunto de ceros de $\varphi_n$ como $Z_{n}$
  \end{enumerate}
\end{definition}

\begin{proposition}\label{prop:singular_values_limit}
  Sea $\mathcal{D}$ el operador multiplicación por $z$ en
  $\ell_2$ y
  $\mathbf{D}=(d_{i,j})_{i,j\ge0}$ su matriz infinita
  . Sea
  $\mathbf{D}_n=(d_{i,j})_{0\le i,j\le n}$ la sección de orden $n+1$. Entonces
  \[
    \lim_{n\to\infty}\|\mathbf{D}_n\|
    \;=\;\|\mathcal{D}\|.
  \]
\end{proposition}

Esta proposición establece una conexión importante entre las normas de las secciones finitas y la norma del operador infinito.

\noindent Utilizando este resultado se tiene que:
$$
  \lim_{n\to \infty} \sigma_{1,n}(\mathbf{D}_n)= \lim_{n \to \infty} \Vert \mathbf{D}_n\Vert= \Vert \mathcal{D}\Vert
$$

\noindent que siempre existe pudiendo ser $\infty$ en el caso no acotado.

Con estos fundamentos de álgebra lineal y teoría de operadores establecidos, podemos ahora adentrarnos en la definición específica de los productos internos de Sobolev.

\section{Productos internos de Sobolev}

\noindent Una vez establecidos los conceptos básicos del álgebra lineal, debemos extender nuestro marco teórico al contexto de espacios funcionales, donde se desarrolla el estudio de los polinomios ortogonales de Sobolev. Para ello, necesitamos introducir los productos de Sobolev, que nos permitirán formalizar los distintos productos internos respecto a los cuales definiremos la ortogonalidad. Estos productos internos pueden tomar diversas formas, ya sean continuos o discretos, reales o complejos, y son fundamentales para caracterizar las propiedades de los polinomios ortogonales asociados.

\begin{definition}\label{def:circum_lebesgue_measure}
  Sea
  \(
  S_{r}(a)\;=\;\bigl\{\,z\in\mathbb{C}\;:\;|z-a|=r\bigr\}
  \)
  la circunferencia de centro \(a\in\mathbb{C}\) y radio \(r>0\).
  Denotaremos por \(dm_{a;r}\) la \textbf{medida de Lebesgue} en \(S_{r}(a)\).
  Para toda función continua \(f:S_{r}(a)\to\mathbb{C}\) definimos
  $$
    \int_{S_{r}(a)} f(z)\,dm_{a;r}(z)
    \;:=\;
    \frac{1}{2\pi}\int_{0}^{2\pi} f\!\bigl(a+re^{i\theta}\bigr)\,d\theta,
  $$
  de modo que \(dm_{a;r}\) queda normalizada
  \(\bigl(dm_{a;r}(S_{r}(a))=1\bigr)\).
\end{definition}

Esta medida es fundamental para definir productos internos en el plano complejo, especialmente en el contexto de los polinomios ortogonales en la circunferencia unitaria o en otras circunferencias.

Los productos internos de Sobolev con los que se trabajará son de la forma:
\begin{enumerate}
  \item Continuo Real con Soportes Compactos
        \begin{equation}
          \langle p(t), q(t) \rangle = \int_{a}^{b} p(t) q(t) w_1(t) \, dt + \int_{c}^{d} p'(t) q'(t) w_2(t) \, dt
        \end{equation}
        donde $w_1(t)$ y $w_2(t)$ son pesos continuos, no negativos y distintos de la función $f \equiv 0$.
  \item Discreto Real con Átomo
        \begin{equation}
          \langle p(t), q(t) \rangle = \int_{a}^{b} p(t) q(t) w_1(t) \, dt + p'(c) q'(c)
        \end{equation}
        con $w_1(t)$ peso continuo, no negativo y distinto de la función $f \equiv 0$.
  \item Continuo Complejo en Circunferencias
        \begin{equation}
          \langle p(z) , q(z) \rangle_{S}\;=\;
          \int_{S_{r_b}(b)}     p(z)\,\overline{q(z)} \, dm_{b;r_b}
          \;+\;
          \int_{S_{r_a}(a)}
          p'(z)\,\overline{q'(z)} \, dm_{a;r_a}.
        \end{equation}
  \item Discreto Complejo en Circunferencia con Átomo
        \begin{equation}
          \langle p(z) , q(z) \rangle_{S}\;=\;
          \int_{S_{r_b}(b)} p(z)\,\overline{q(z)}\,dm_{b;r_b}
          \;+\;
          p'(c)\,\overline{q'(c)}.
        \end{equation}
\end{enumerate}

Mas generalmente,
\begin{definition}\label{def:sobolev_product}
  \cite{escribano2025boundedness}
  Sea \(\boldsymbol{\mu}=(\mu_1,\mu_2)\) un conjunto de medidas de Borel,
  positivas y finitas, con soportes \(\operatorname{supp}(\mu_1),\operatorname{supp}(\mu_2)\subset\mathbb C\) y \( \mu_1\) con soporte infinito,
  y sea \(p,q\in\mathbb P[z]\).
  Definimos el \emph{producto de Sobolev} asociado a
  \(\boldsymbol{\mu}\) como
  \[
    \langle p,q \rangle_{S}
    \;=\;
    \int_{\mathbb C} p(z)\,\overline{q(z)}\,\mathrm d\mu_{1}(z) + \int_{\mathbb C} p'(z)\,\overline{q'(z)}\,\mathrm d\mu_{2}(z).
  \]
\end{definition}

Para los casos discretos, la medida de Dirac juega un papel crucial.

\begin{definition}\label{def:dirac_measure}
  Sea $a\in\mathbb{C}$.
  La \emph{medida de Dirac} en $a$, denotada $\delta_{a}$, es la medida de Borel definida por
  \[
    \delta_{a}(A)\;=\;
    \begin{cases}
      1, & a\in A,    \\
      0, & a\notin A,
    \end{cases}
    \qquad A\subset\mathbb C \text{ Borel}.
  \]

  Para cualquier función medible \(f\colon\mathbb{C}\to\mathbb{C}\) se cumple
  \[
    \int_{\mathbb C} f(z)\,\mathrm d\delta_{a}(z)=f(a).
  \]

  En particular, si \(p\) es un polinomio,
  \[
    |p(a)|^{2}\;=\;
    \int_{\mathbb C} |p(z)|^{2}\,\mathrm d\delta_{a}(z).
  \]

  Cuando se desea ponderar la masa, se emplea el “átomo ponderado”
  \(
  \lambda\,\delta_{a}\;( \lambda>0),
  \)
  para el cual
  \(
  \displaystyle \int f\,\mathrm d(\lambda\delta_{a})=\lambda\,f(a).
  \)
\end{definition}

\paragraph{Observación:}
En los casos discretos con atomo, la medida utilizada es la medida de Dirac, es decir, se considera un atomo en un punto específico del espacio complejo, lo que permite evaluar la función en ese punto y ponderar su contribución al producto interno.

\begin{definition}
  \label{rem:R_support}
  \noindent En un espacio de Sobolev asociado a dos medidas se define el número  $R$
  $$
    R= \max \{ \vert z \vert : z\in {\it supp}(\mu_1)\cup {\it supp}(\mu_2)\},
  $$
\end{definition}
\noindent En particular, en los casos tratados:
\begin{enumerate}
  \item Caso continuo real
        $$
          R= \max \{ \vert a \vert,\vert b \vert, \vert c \vert, \vert d \vert\}
        $$
  \item Caso continuo con átomo:
        $$
          R= \max \{\vert a \vert, \vert b \vert, \vert c \vert\}
        $$
  \item Caso continuo complejo
        $$
          R=\max \{ z \in S_{r_a}(a)\cup S_{r_b}(b)\}= \max \{ \vert a \vert + r_{a}, \vert b \vert + r_b \}
        $$
  \item Caso complejo con átomo
        $$
          R= \max\{ \max{\vert z \vert : z\in S_{r_a}(a)}, \vert c \vert \}= \max \{ \vert a \vert + r_a,\vert c \vert \}
        $$
\end{enumerate}

\noindent Hay que destacar que este número, en el caso de un productos internos asociados a una sóla medida, es decir, el caso standard de polinomios ortogonales asociados a una medida del tipo
$$
  <p,q>=\int p(z) \overline{q(z)} \; d\mu
$$

\noindent
coincide con
$$
  R=\max \{ \vert z \vert : z\in {\it supp}(\mu)\}
$$

\noindent Es conocido que cuando la medida tiene soporte acotado,
$$\lim_{n\to \infty} \sigma_{1,n} = \Vert \mathcal{D}\Vert = R$$


Este valor $R$ proporciona una medida del tamaño de la región donde se definen las medidas del producto de Sobolev. Al igual que con los productos internos generales, los productos de Sobolev también admiten una representación matricial.

\begin{definition}\label{def:sobolev_matrix}
  \cite{escribano2025boundedness}
  Sea $\boldsymbol{\mu}=(\mu_{1},\mu_{2})$ un conjunto de medidas de Borel,
  positivas, con soporte compacto en $\mathbb{C}$.
  Denotemos por
  \[
    \textbf{M}_j:=\textbf{M}(\mu_j)=\bigl(c^{(j)}_{n,m}\bigr)_{n,m\ge0},
    \qquad
    c^{(j)}_{n,m}:=\int_{\mathbb C} z^{\,n}\,\overline{z}^{\,m}\,{\rm d}\mu_j(z),
    \quad j=1,2,
  \]
  las respectivas matrices de momentos (cada una HSPD).
  \vspace{0.4em}
  La \emph{matriz de momentos Sobolev}
  asociada al sistema \(\mu\) es entonces
  \[
    \textbf{M}_S \;=\;
    \textbf{M}_1 + \boldsymbol{\Lambda}\,\textbf{M}_2\,
    \boldsymbol{\Lambda}^{\,*}
  \]

  donde $\boldsymbol{\Lambda} = (\lambda_{i,j})_{i,j\ge0}$ es la matriz infinita que representa el operador derivación definida por:
  \[
    \lambda_{i,j} = \begin{cases}
      j & \text{si } i = j-1  \\
      0 & \text{en otro caso}
    \end{cases}
  \]
  o equivalentemente:
  \[
    \boldsymbol{\Lambda} = \begin{pmatrix}
      0      & 0      & 0      & 0      & 0      & \cdots \\
      1      & 0      & 0      & 0      & 0      & \cdots \\
      0      & 2      & 0      & 0      & 0      & \cdots \\
      0      & 0      & 3      & 0      & 0      & \cdots \\
      \vdots & \vdots & \vdots & \vdots & \vdots & \ddots
    \end{pmatrix}
  \]

  Luego el \emph{producto matricial interno de Sobolev} asociada a un conjunto de matrices HPSD \(\{\textbf{M}_1,\textbf{M}_2\}\) con $\textbf{M}_1$ siendo HPD en $\mathbb{P}[z]$ se define como:

  $$
    \langle p(t), q(t) \rangle_{\textbf{M}_{S}} = v \Bigl( \textbf{M}_1 + \boldsymbol{\Lambda}\,\textbf{M}_2\,
    \boldsymbol{\Lambda}^{\,*}\Bigr) w^{*}.
  $$
\end{definition}

La matriz $\boldsymbol{\Lambda}$ representa el operador de derivación en la base monomial. Cuando se utilizan medidas de Dirac, la matriz de momentos asociada tiene una estructura particular.

\begin{definition}\label{def:dirac_measure_matrix}
  \cite{escribano2025boundedness}
  Sea \(\delta_{a}\) la medida de Dirac concentrada en \(a\in\mathbb C\).
  Sus momentos son
  \(
  c_{n,m}= \displaystyle\int_{\mathbb C} z^{\,n}\,\overline{z}^{\,m}\,d\delta_{a}(z)
  = a^{\,n}\,\overline{a}^{\,m},
  \quad n,m\ge0.
  \)

  La \emph{matriz de momentos} asociada,
  denotada \(M(a)=M(\delta_{a})\), es por tanto
  \[
    M(a)=\bigl(a^{\,n}\,\overline{a}^{\,m}\bigr)_{n,m\ge0}
    =
    \begin{pmatrix}
      1      & \overline{a}      & \overline{a}^{2}      & \overline{a}^{3}      & \dots  \\
      a      & |a|^{2}           & a\overline{a}^{2}     & a\overline{a}^{3}     & \dots  \\
      a^{2}  & a^{2}\overline{a} & |a|^{4}               & a^{2}\overline{a}^{3} & \dots  \\
      a^{3}  & a^{3}\overline{a} & a^{3}\overline{a}^{2} & |a|^{6}               & \dots  \\
      \vdots & \vdots            & \vdots                & \vdots                & \ddots
    \end{pmatrix}.
  \]

  Nótese que dicha matriz de momentos no define un producto interno.
\end{definition}

Aunque $M(a)$ por sí sola no define un producto interno (ya que no es definida positiva, solo semidefinida positiva), puede ser un componente en la construcción de productos internos de Sobolev más generales.

Escribano y Gonzalo generalizaron la definición de producto interno de Sobolev a producto interno de Sobolev matricial, considerando matrices que no necesariamente provienen de medidas de Borel.

\begin{definition}\label{def:sobolev_matrix_product}
  \cite{escribano2025boundedness} Sea $\{\textbf{M}_1,\textbf{M}_2\}$ un conjunto de matrices HPSD, siendo $\textbf{M}_1$ HPD, el producto matricial de Sobolev inducido, denotado por $\langle \cdot, \cdot \rangle_{\textbf{M}_{\textbf{S}}}$, se define como:
  \begin{equation}
    \langle p(z), q(z) \rangle_{\textbf{M}_{\textbf{S}}} = \langle p(z), q(z) \rangle_{\textbf{M}_1} + \langle p'(z), q'(z) \rangle_{\textbf{M}_2},\quad \forall p(z), q(z) \in \mathbb{P}[z]
  \end{equation}
\end{definition}

\paragraph{Observación}
Cuando \(\mathbf{M}_{1}\) es la matriz de momentos de una medida \(\mu_{1}\) y
\(\mathbf{M}_{2}\) de \(\mu_{2}\), los cuatro casos (1.3)-(1.6) se recuperan como casos particulares.

Con estas definiciones establecidas, podemos formular el problema central de este trabajo.

\section{Planteamiento del problema}

El problema central de investigación que motiva este trabajo es el siguiente: Sean $\mu_1$ y $\mu_2$ dos medidas de Borel positivas con soporte compacto, con $\mu_1$ con soporte infinito, sea $\langle\cdot,\cdot\rangle_S$ el producto interno de Sobolev asociado y sea $Z$ el conjunto de ceros de los polinomios ortonormales de Sobolev $\{\varphi_i\}_{i\ge 0}$. ¿Es $Z$ acotado?

Esta cuestión fundamental no ha sido completamente resuelta en la literatura, especialmente en el caso complejo cuando los soportes son disjuntos o cuando no se cumple la hipótesis de dominación secuencial, conceptos que se verán después.

El estudio de la acotación de los ceros de los polinomios ortogonales de Sobolev se
traslada, mediante el uso de la representación matricial del
operador multiplicación $\mathcal{D}$, al análisis del \emph{radio espectral} de las secciones de la matriz infinita. Sabemos que si la sucesión \(\{\rho(\mathbf{D}_{n})\}_{n\ge0}\) esta acotada, entonces el \emph{conjunto de ceros}
\(Z\) estará acotado.

\textbf{Objetivo: }
Determinar condiciones necesarias y suficientes para la acotación del \emph{conjunto de ceros} en cada uno de los casos (1.3)-(1.6) de productos internos de Sobolev, estableciendo relaciones entre:
\begin{itemize}
  \item la acotación de $Z$,
  \item la acotación del operador multiplicación $\mathcal{D}$,
  \item el comportamiento asintótico de $\rho(\mathbf{D}_{n})$.
  \item el comportamiento asintótico $\sigma_{1,n}$ y $\sigma_{2,n}$.
\end{itemize}

\paragraph{Metodología: }
Para la parte computacional emplearemos
\textsc{Maple}\, con el fin de:
\begin{enumerate}
  \item generar las matrices \(\mathbf{D}_{n}\) para valores crecientes de \(n\),
  \item calcular los radios espectrales $\rho(\mathbf{D}_{n})$ y analizar su comportamiento asintótico,
  \item visualizar los polinomios ortogonales y localizar sus ceros en el plano complejo,
  \item contrastar numéricamente las cotas teóricas obtenidas y explorar casos donde las condiciones clásicas de dominación secuencial no se satisfacen.
\end{enumerate}

Antes de abordar este problema, es importante revisar el estado actual del conocimiento en esta área.

\section{Estado del arte}

En la literatura se ha abordado el problema de los ceros de polinomios ortogonales de Sobolev mayoritariamente desde la perspectiva de la localización general de los ceros.

Primero vamos a probar que la acotacion del operador multiplicación $\mathcal{D}$ en $\mathbb{P}[z]$ es condicion suficiente para la acotación de los ceros de los polinomios ortogonales respecto a un producto inducido por una matriz HPD.

\begin{proposition}\label{prop:bounded_zeros}
  \cite{escribano2025boundedness} Sea \textbf{M} una matriz HPD y $\{\varphi_n(z)\}_{n=0}^{\infty}$ la secuencia de polinomios ortonormales asociados al producto interno inducido por \textbf{M}. Si $\mathcal{D}$ está acotado en $\mathbb{P}[z]$, el conjunto de ceros de los polonomios ortonormales asociados a \textbf{M} está acotado. Además,
  $$Z=\{z\in\mathbb{C}: \exists n \in \mathbb{N},\varphi_n(z)=0\} \subset \{z\in\mathbb{C}:|z|<\|\mathcal{D}\|\}.$$
\end{proposition}

\begin{proof}
  Asumamos que para $z_0 \in Z$ existe un $n \in \mathbb{N}$ tal que $\varphi_n(z_0)=0$. Luego, $\varphi_n(z)=(z-z_0)P_{n-1}(z)$, $0 \neq P_{n-1} \in \mathbb{P}_{n-1}[z]$. Por lo tanto, $\varphi_n(z_0)=zP_{n-1}(z)-z_0P_{n-1}(z)$ y porque $\varphi_n(z)$ y $P_{n-1}(z)$ son ortogonales en \textit{M} se sigue que
  $$\|zP_{n-1}(z)\|^2=\langle\varphi_n(z)+z_0P_{n-1}(z),\varphi_n(z)+z_0P_{n-1}(z)\rangle = \|\varphi_n(z)\|^2+|z_0|^2\|P_{n-1}(z)\|^2,$$
  por lo que,
  $$|z_0|^2\|P_{n-1}(z)\|^2=\|zP_{n-1}(z)\|^2- \|\varphi_n(z)\|^2<\|zP_{n-1}(z)\|^2=$$
  $$\|zP_{n-1}(z)\|^2=\|\mathcal{D}(P_{n-1})\|^2\leq\|\mathcal{D}\|^2\|P_{n-1}(z)\|^2.$$
  Como $\|P_{n-1}(z)\| \neq 0$, se concluye que
  $$|z_0|^2<\|\mathcal{D}\|,$$
  como se requería.
\end{proof}

Este resultado fundamental y usado extensamente en la literatura establece que los ceros de polinomios ortogonales de Sobolev están contenidos en un disco centrado en cero cuyo radio es la norma del operador $\mathcal{D}$, si dicho operador está acotado. Veamos resultados que nos permiten acotar la norma del operador multiplicación $\mathcal{D}$.

Para acotar el operador multiplicación, G. López y H. Pijeira \cite{lagomasino1999zero} introdujeron el concepto de medidas secuencialmente dominadas.

\begin{definition}
  \cite{lagomasino1999zero} Diremos que el producto de Sobolev con dos medidas es secuencialmente dominado si
  $$supp(\mu_2) \subset supp(\mu_{1}) , \qquad j = 1,2.$$ y
  $$d\mu_2 = f_{1}d\mu_{1}, \qquad f_{1} \in L_{\infty}(\mu_{1}).$$
\end{definition}

Esta condición de dominación secuencial es clave para garantizar la acotación del operador multiplicación en ciertos casos.

En el caso continuo real, se tiene que como consecuencia de la Definición \ref{def:seq_dominated}:

\begin{corollary}
  El producto interno de Sobolev
  $$\langle p(t), q(t) \rangle = \int_{a}^{b} p(t) q(t) w_1(t) \, dt + \int_{c}^{d} p'(t) q'(t) w_2(t) \, dt$$
  es secuencialmente dominado cuando se cumple que $[c,d] \subset [a,b]$ y $\exists c > 0, w_2(t)\leq c \cdot w_1(t)$ continuos.
\end{corollary}

Para trabajar con matrices de momentos introduciremos la definición para esta noción de secuencial dominacion y un teorema análogo que nos dará la acotación. Estos conceptos fueron generalizados por Escribano y Gonzalo \cite{escribano2025boundedness} para el caso matricial.

\begin{definition}\label{def:seq_dominated}
  \cite{escribano2025boundedness} Sea $\{\textbf{M}_1,\textbf{M}_2\}$ un conjunto de matrices HSPD. Diremos que $\{\textbf{M}_1,\textbf{M}_2\}$ es secuencialmente dominada si existe una constante positiva $C$ tal que
  $$\textbf{M}_2 \leq C \cdot \textbf{M}_{1}.$$
  donde $\textbf{A} \leq \textbf{B}$ significa que $v\textbf{A}v^* \leq v\textbf{B}v^*$ para todo $v \in c_{00}$.
\end{definition}

Como consecuencia de \cite{escribano2025zeros} Proposición 5, se tiene que el producto Sobolev
$$\langle p(z) , q(z) \rangle_{S}\;=\;
  \int_{S_{r_b}(b)} p(z)\,\overline{q(z)} \, dm_{b;r_b}
  \;+\;
  \int_{S_{r_a}(a)} p'(z)\,\overline{q'(z)} \, dm_{a;r_a}.$$
es matricialmente secuencialmente dominado, es decir, $\{\textbf{M}(m_{b;r_b}),\textbf{M}(m_{a;r_a})\}$ forman un conjunto de matrices secuencialmente dominado cuando se cumple que $S_{r_a}(a)\subset \mathbb{D}_{r_b}(b)$, donde $\mathbb{D}_{r_b}(b) = \{z \in \mathbb{C} : |z - b| < r_b\}$ denota el disco abierto de centro $b$ y radio $r_b$.

\begin{theorem}\label{thm:bounded_operator}
  \cite{escribano2025boundedness} Sea $\{\textbf{M}_1,\textbf{M}_2\}$ un conjunto de matrices HPD secuencialmente dominado. Asuma que para $j=1,2$ el operador multiplicación $\mathcal{D}_j$ está acotado en $(\mathbb{P}[z],\langle \cdot,\cdot \rangle)_{\textbf{M}_j}$. Considera el producto matricial de Sobolev siguiente:
  $$\langle p(t),q(t) \rangle_{\textbf{M}_S}=\langle p(t),q(t) \rangle_{\textbf{M}_1}+\langle p'(t),q'(t) \rangle_{\textbf{M}_2}.$$
  Entonces, el operador multiplicación $\mathcal{D}$ está acotado en $(\mathbb{P}[z],\langle \cdot,\cdot \rangle)_{\textbf{M}_S}.$
\end{theorem}

Este teorema es una herramienta poderosa para establecer la acotación del operador multiplicación en el contexto de productos de Sobolev matriciales. Un corolario directo se aplica a los productos de Sobolev definidos en circunferencias.

\begin{corollary}\label{cor:seq_dominated_sobolev}
  Como consecuencia del Teorema \ref{thm:bounded_operator}, para el producto interno de Sobolev
  $$\langle p(z) , q(z) \rangle_{S}\;=\;
    \int_{S_{r_b}(b)} p(z)\,\overline{q(z)} \, dm_{b;r_b}
    \;+\;
    \int_{S_{r_a}(a)} p'(z)\,\overline{q'(z)} \, dm_{a;r_a}.$$
  el operador multiplicación $\mathcal{D}$ está acotado en $(\mathbb{P}[z],\langle \cdot,\cdot \rangle)_{S}$ cuando se cumple que $S_{r_a}(a)\subset \mathbb{D}_{r_b}(b)$. Y por tanto el conjunto $Z$ está acotado.
\end{corollary}

Sin embargo, como consecuencia de \cite{escribano2025zeros} Proposición 6, podemos asegurar que para el producto Sobolev
$$
  \langle p(z) , q(z) \rangle_{S}\;=\;
  \int_{S_{r_b}(b)} p(z)\,\overline{q(z)} \, dm_{b;r_b}
  \;+\;
  \int_{S_{r_a}(a)} p'(z)\,\overline{q'(z)} \, dm_{a;r_a}.
$$
El operador multiplicación $\mathcal{D}$ está acotado en $(\mathbb{P}[z],\langle \cdot,\cdot \rangle)_{S}$ cuando se cumple que $a \in \mathbb{D}_{r_b}(b)$ y $r_a > 0$, incluso cuando no estamos en un caso matricialmente secuencialmente dominado. Este resultado es particularmente interesante porque muestra que hay casos donde la acotación del operador puede obtenerse bajo condiciones geométricas más generales que la dominación secuencial.

Para extender estos resultados a productos de Sobolev que incorporan términos discretos, como la evaluación de la derivada en un punto específico, necesitamos introducir el concepto de punto de evaluación acotado.

Los puntos de evaluación acotados constituyen una herramienta fundamental para extender los resultados de acotación del operador multiplicación a productos internos de Sobolev con componentes discretas. Como veremos en los resultados siguientes, cuando un punto es de evaluación acotada respecto a la medida principal $\mu_1$, es posible garantizar la acotación del operador multiplicación en productos internos de Sobolev que incluyen términos discretos como $|p'(a)|^2$, lo que a su vez asegura la acotación del conjunto de ceros de los polinomios ortogonales asociados.

\begin{definition}\label{def:bounded_eval_point}
  \cite{escribano2025zeros}
  Decimos que \(a\in\mathbb{C}\) es un \textbf{punto de evaluación acotado} de una medida \(\mu\) (abreviado \textbf{bpe} de \(\mu\)) si existe una constante \(C>0\) tal que
  \[
    |p(a)|^{2} \leq C \int |p(z)|^{2}\,d\mu(z), \qquad p(z)\in\mathbb{P}[z].
  \]
\end{definition}

\begin{definition}\label{def:poly_convex_hull}
  \cite{escribano2025zeros}
  Para un conjunto compacto \(K \subset \mathbb{C}\), la \textbf{envoltura convexa polinómica} de \(K\), denotada por \(P_C(K)\), se define como
  \[
    P_C(K) =
    \left\{ z \in \mathbb{C} :\, |p(z)| \leq \max_{\xi\in K}|p(\xi)|\ \text{ para todo } p\in\mathbb{P}[z] \right\}.
  \]
\end{definition}

La envoltura convexa polinómica juega un papel importante en la caracterización de los puntos de evaluación acotados.

En particular, por \cite{Conway1978}, \cite{Gamelin1969}
\begin{enumerate}
  \item Para el caso real, cuando $K$ es union de intervalos, $P_C(K)$ coincide con $K$.
  \item Para el caso en el que $K$ es una circunferencia, cuando $K$ es $S_{r}(a)$, $P_C(K)$ coincide con $\mathbb{D}_{r}(a)$.
\end{enumerate}

Estos resultados conectan la geometría del soporte de la medida con las propiedades analíticas de los puntos de evaluación.

\begin{lemma}\label{lem:bounded_eval_point_in_hull}
  \cite{escribano2025zeros} Todo \textbf{punto de evaluación acotado} de una medida $\mu$ está contenido en la envoltura polinómica convexa del soporte de $\mu$.

  \noindent Además, todo \textbf{punto de evaluación acotado} de una medida $\mu$ es un punto de evaluación de la matriz de momentos asociada a $\mu$.
\end{lemma}

La noción de punto de evaluación acotado se extiende naturalmente al contexto de matrices de momentos.

\begin{definition}\label{def:bounded_eval_point_matrix}
  \cite{escribano2025zeros} Sea $M$ una matriz hermítica positiva definida (HPD) y $a\in\mathbb{C}$.
  Decimos que $a$ es un \textbf{punto de evaluación acotado} (bpe) de $M$ si existe una constante $C>0$ tal que
  \[ |p(a)|^{2}\;\le\;C\,\|p(z)\|_{M}^{2}, \quad p(z)\in\mathbb{P}[z]. \]
\end{definition}

Existe una equivalencia útil entre ser un punto de evaluación acotado y la dominación secuencial.

\begin{proposition}\label{prop:equiv_bounded_eval_point}
  \cite{escribano2025zeros} Sea $M$ una matriz HPD y $a\in\mathbb{C}$. Entonces son equivalentes:
  \begin{enumerate}
    \item $a$ es un \textbf{punto de evaluación acotado} de $M$;
    \item el conjunto $\{M,\,M(a)\}$ es secuencialmente dominado;
  \end{enumerate}
\end{proposition}

Esta equivalencia permite aplicar resultados de dominación secuencial al estudio de puntos de evaluación acotados. La siguiente proposición muestra cómo la propiedad de ser un punto de evaluación acotado se preserva bajo la adición de un término de derivada.

\begin{proposition}\label{prop:bounded_derivative_eval}
  \cite{escribano2025zeros} Sea $M$ una matriz HPD tal que el operador de multiplicación $D$
  es acotado en $\bigl(\mathbb{P}[z],\|\cdot\|_{M}\bigr)$, y sea $a\in\mathbb{C}$ punto de evaluación acotado de $M$.
  Entonces $D$ también es acotado en $\mathbb{P}[z]$ respecto a la norma
  \[
    \|p(z)\|_{\textbf{M}_{S}}^{2}\;=\;\|p(z)\|_{M}^{2}\;+\;
    \bigl|p'(a)\bigr|^{2}.
  \]
\end{proposition}

Un corolario importante de este resultado se aplica a productos de Sobolev discretos en circunferencias.

\begin{corollary}\label{cor:bounded_zeros_circle}
  \cite{escribano2025zeros}
  Sea $m_{b;r_b}$ una medida de Lebesgue soportada en una circunferencia $S_{r_b}(b)$.
  Supongamos que $a\in \mathbb{D}_{r_b}(b)$ y consideremos la norma de Sobolev discreta
  \[
    \|p\|_{S}^{2}= \int_{S_{r_b}(b)} |p(z)|^{2}\,dm_{b;r_b}\;+\;\bigl|p'(a)\bigr|^{2},
    \qquad p\in\mathbb{P}[z].
  \]
  Entonces el operador multiplicación $\mathcal{D}\colon p(z)\mapsto z\,p(z)$
  es acotado en $\bigl(\mathbb{P}[z],\|\cdot\|_{S}\bigr)$.
  En consecuencia, el conjunto de ceros de los polinomios ortogonales de Sobolev asociados está acotado.
\end{corollary}

Estos resultados proporcionan condiciones suficientes para la acotación de los ceros en varios escenarios importantes.

Además de las técnicas basadas en matrices de momentos y acotación de operadores, existen otros enfoques en la literatura para estudiar la acotación del conjunto de ceros de polinomios ortogonales de Sobolev. Estos métodos alternativos no requieren necesariamente el uso de representaciones matriciales ni el análisis de la acotación del operador multiplicación.

Un enfoque geométrico importante se basa en el análisis directo de las propiedades de los soportes de las medidas involucradas. Por ejemplo, cuando los soportes de las medidas $\mu_1$ y $\mu_2$ son disjuntos, es posible obtener cotas explícitas para la localización de los ceros utilizando técnicas de análisis complejo y teoría de potencial.

Otro método relevante utiliza las propiedades asintóticas de los coeficientes de las relaciones de recurrencia de tres términos que satisfacen los polinomios ortogonales de Sobolev. Este enfoque, desarrollado por varios autores, permite estudiar el comportamiento de los ceros sin necesidad de construir explícitamente las matrices asociadas.

También se han desarrollado técnicas basadas en desigualdades integrales y métodos variacionales que proporcionan información sobre la distribución asintótica de los ceros. Estos métodos son particularmente útiles cuando se busca información cualitativa sobre el comportamiento de los ceros más que cotas cuantitativas precisas.

Un ejemplo concreto de estas técnicas alternativas se presenta en el siguiente resultado para soportes disjuntos:

\begin{theorem}\label{thm:disjoint_support}
  \cite{lagomasino2001sobolev} En un producto Sobolev del Caso Real Continuo, si $\operatorname{co}(\operatorname{supp}(\mu_1)) \cap \operatorname{co}(\operatorname{supp}(\mu_2)) = \emptyset$, entonces, los ceros de $\varphi_n$ están contenidos en el disco centrado en el punto de $\operatorname{co}(\operatorname{supp}(\mu_2))$ más alejado de $\operatorname{co}(\operatorname{supp}(\mu_1))$ y
  radio el diámetro de $\operatorname{co}(\operatorname{supp}(\mu_1)) \cup \operatorname{co}(\operatorname{supp}(\mu_2))$. Donde $\operatorname{co}(K)$ es la envoltura convexa de un conjunto \(K\),
  \[
    \operatorname{co}(K)=\Bigl\{\sum_{j=1}^m \lambda_j x_j \;:\;
    x_j\in K,\;
    \lambda_j\ge 0,\;
    \sum_{j=1}^m \lambda_j =1,\;
    m\in \mathbb{N}\Bigr\};
  \]
  es decir, el menor conjunto convexo que contiene a \(K\) o, equivalentemente, la intersección de todos los subconjuntos convexos cerrados de \(C\) que contienen a \(K\).
\end{theorem}